\section*{Глава 1. Анализ предметной области}
\addcontentsline{toc}{section}{Глава 1. Анализ предметной области}
\stepcounter{section}

\subsection{Обзор существующих платформ электронной коммерции}

Разработка платформы электронной коммерции по модели маркетплейса для продажи цифровой техники является сложной инженерной задачей. Специфика 
предметной области требует обработки глубоко вложенных характеристик товаров (для реализации системы сложных фильтров), обеспечения 
строгой изоляции данных независимых продавцов и высокой скорости отклика пользовательского интерфейса.

Анализ современного рынка готового программного обеспечения выявляет три основных подхода к созданию подобных платформ. Рассмотрим 
их технические особенности:

\begin{enumerate}
    \item \textbf{Сборки на базе WordPress (WooCommerce).} Open Source\footnotemark решение. Главный недостаток — использование 
	потенциально большого количества атрибутов у сущностей в базе данных. При работе с большим каталогом электроники и множественными 
	фильтрами производительность системы падает.
    
    \item \textbf{Специализированные CMS (CS-Cart Multi-Vendor).} Коробочное решение с готовой логикой для маркетплейсов. Недостаток — устаревающая 
	монолитная архитектура с серверным рендерингом, что приводит к постоянным перезагрузкам страниц и снижению отзывчивости интерфейса.
    
    \item \textbf{Корпоративные решения (1C-Bitrix Enterprise).} Мощная платформа с интеграцией складского учета. Отличается высокой стоимостью 
	лицензии и закрытостью ядра, что сильно усложняет и удорожает внедрение современных одностраничных приложений (SPA) на базе API.
\end{enumerate}

\footnotetext{Распространяется под свободной лицензией GNU GPL.}

Для структурирования полученных данных был проведен сравнительный анализ по ключевым критериям разработки. Результаты представлены 
в таблице~\ref{tab:competitors}.

\begin{table}[H]
	\centering
	\caption{Сравнительный анализ платформ электронной коммерции}
	\label{tab:competitors}
	\begin{tabular}{|p{0.25\textwidth}|p{0.22\textwidth}|p{0.18\textwidth}|p{0.25\textwidth}|}
		\hline
		\textbf{Название} & \textbf{Тип \linebreak архитектуры} & \textbf{Гибкость \linebreak модели БД} & \textbf{Готовность \linebreak к SPA/API} \\ \hline
		WordPress + \linebreak WooCommerce & Монолит & Низкая & Низкая \\ \hline
		CS-Cart \linebreak Multi-Vendor & Монолит & Средняя & Средняя \\ \hline
		1C-Bitrix \linebreak Enterprise & Монолит + REST API & Высокая & Требует написания \linebreak прослойки API \\ \hline
	\end{tabular}
\end{table}

Детальный анализ существующих программных решений показал, что подавляющее большинство из них опирается на монолитную архитектуру 
и устаревшие подходы к рендерингу, что делает их медленными при работе с большими каталогами электроники. 
Использование готовых CMS накладывает ограничения на структуру базы данных, усложняя реализацию быстрых фильтров.

В связи с этим подтверждается целесообразность разработки собственного программного продукта. Разделение системы на независимый серверный API 
и клиентское приложение позволит добиться максимальной гибкости ролевой модели, высокой скорости отклика интерфейса и 
оптимизации работы с реляционной базой данных.

\subsection{Анализ программных инструментов разработки}

Для реализации проекта выбран современный стек технологий, обеспечивающий разделение приложения на независимые клиентскую и серверную части. Выбор инструментов определен требованиями 
к масштабируемости, скорости отклика интерфейса и безопасности данных:

\begin{itemize}
    \item \textbf{Frontend-библиотека (React.js):} На современном рынке клиентской разработки доминируют такие фреймворки, как Angular, Vue.js и библиотека React. Angular 
	обладает излишне жесткой архитектурой, которая избыточна для данного проекта, а Vue.js имеет менее обширную систему готовых инструментов для управления состояниями. Выбор библиотеки 
	React обусловлен компонентным подходом, использованием виртуальной древовидной структуры объектов, что обеспечивает высокую 
	скорость перерисовки интерфейса каталога при применении множества фильтров, а также наличием официальных рекомендаций по разработке интерфейсов~\cite{react_docs}.
    \item \textbf{Стилизация (Tailwind CSS):} Для стилизации компонентов традиционно применяются CSS-препроцессоры или готовые UI-библиотеки, например, Bootstrap. Стоит отметить, что 
	препроцессоры усложняют поддержку кода при росте проекта, а UI-библиотеки сильно ограничивают уникальность дизайна. Фреймворк Tailwind CSS был выбран для ускорения верстки интерфейса и минимизации 
	объема CSS-кода за счет использования утилитарных классов, подробное описание которых приводится в документации~\cite{tailwind_docs}.
    \item \textbf{Backend-платформа (Python + FastAPI):} При выборе серверной технологии рассматривались такие популярные решения, как Node.js и Python-фреймворки (Django, Flask). Node.js требует сторонних 
	решений для строгой валидации данных, Django является тяжеловесным монолитом, а Flask не имеет встроенной асинхронности. В свою очередь, FastAPI обеспечивает высокую производительность за счет асинхронности и 
	автоматическую генерацию документации с помощью Swagger, что необходимо для обработки запросов от множества пользователей одновременно~\cite{fastapi_docs}.
    \item \textbf{СУБД (PostgreSQL):} Для хранения данных рассматривались нереляционные (MongoDB) и реляционные (MySQL, PostgreSQL) СУБД. Нереляционные базы не подходят для платформ электронной коммерции, так как не 
	обеспечивают надежность финансовых транзакций. В сравнении с MySQL, выбор PostgreSQL аргументирован лучшей поддержкой сложных связей и наличием типа данных JSONB, что позволит эффективно хранить динамические 
	характеристики электронной техники при сохранении высокой надежности реляционной структуры~\cite{postgresql_docs}.
\end{itemize}

Для реализации дополнительного функционала клиентской части используется библиотека \texttt{react-router-dom}, обеспечивающая маршрутизацию страниц без перезагрузки браузера.

В серверной части приложения для взаимодействия с СУБД PostgreSQL применяется ORM-библиотека SQLAlchemy. Она позволяет абстрагироваться 
от прямых SQL-запросов, предотвращает уязвимости к SQL-инъекциям и эффективно работает с объектами базы данных на языке Python. Для обеспечения безопасности доступа пользователей используется протокол JSON Web Token. Его 
выбор обусловлен раздельной архитектурой веб-приложения: токены не требуют хранения сессий на сервере и позволяют безопасно передавать данные о ролях пользователей между независимыми клиентом и сервером.

\subsection{Формулировка цели и задач работы}

На основе проведенного анализа инструментов и выявленных проблем существующих маркетплейсов была сформулирована цель курсового проекта.

\textbf{Цель работы} --- проектирование и разработка веб-платформы маркетплейса электронной техники, обеспечивающей взаимодействие 
продавцов и покупателей с системой аналитической отчетности.


Для достижения поставленной цели необходимо решить следующие \textbf{задачи}:
\begin{enumerate}
    \item Провести анализ предметной области и обосновать выбор программных инструментов разработки;
    \item Сформировать функциональные требования к платформе и описать бизнес-логику основных ролей пользователей;
    \item Спроектировать модель данных и макеты пользовательских интерфейсов;
    \item Разработать клиентскую часть приложения (карточки товаров, корзина и личные кабинеты);
    \item Реализовать серверную часть с модулями аутентификации, управления каталогом и обработки заказов;
    \item Выполнить развертывание разработанной системы на хостинге и подготовить отчетную документацию.
\end{enumerate}